\section{Exact optimization over an affine family of cyclic monodromies}
\label{sec:affine-cyclic-family}

The input in this section is an explicitly supplied affine modular family.
Let $m\geq 1$ be encoded in binary, let $R=\mathbb Z/m\mathbb Z$, and let
\begin{equation}
 S=\{z\in R^n:Az=b\},\qquad h(z)=c+Hz\in R^\beta,
 \label{eq:affine-seam-family}
\end{equation}
where $A\in R^{q\times n}$ and $H\in R^{\beta\times n}$.
We charge the listed variable dimension $n$ as part of the input size,
even if every matrix has zero rows. More precisely, the size parameter is
$n+q+\beta$ plus the binary length of $m$ and all listed coefficients.
This is the usual explicit-variable model, and it is also the appropriate
output-sensitive measure for returning an $n$-entry witness. A short JSON
integer declaring an enormous number of unused variables is not an
encoding of those variables in this size convention. Initial coefficients
are read and reduced modulo $m$ before the bounded-residue arithmetic.
For a connected base with translation monodromy generated by the entries of
$h(z)$, the number of components is
\begin{equation}
 \kappa(z)=\gcd(m,h_1(z),\ldots,h_\beta(z)).
 \label{eq:affine-component-objective}
\end{equation}
Indeed, every orbit is a coset of the subgroup generated by these translations,
which is the subgroup of multiples of the displayed gcd.  Representatives of
the entries do not affect this gcd.  With $\beta=0$ the convention is
$\gcd(m)=m$, as required for the trivial action on $m$ sheets.

This is an optimization query about a supplied family, not a classification
of all its marked members.  Distinct assignments can have the same component
count while retaining different attachment data.  In particular, the theorem
below does not eliminate an $m^\beta$ lower bound for explicitly listing that
many inequivalent marked states.  When a pre-existing gluing model requires
every block cover to be connected, global connectedness can already be forced
by the assembly graph; the nontrivial use here allows disconnected local
patches, or other affine families in which the monodromy subgroup varies.

\subsection{An affine coset contains a vector attaining its content ideal}

For a vector $v\in R^\beta$ and generator columns $u_1,\ldots,u_k\in R^\beta$,
define
\[
 \gamma=\gcd\bigl(m,\text{all entries of }v,u_1,\ldots,u_k\bigr).
\]
All gcds below are positive divisors of $m$.

\begin{theorem}[Attained affine content]
\label{thm:attained-affine-content}
There exist $t_1,\ldots,t_k\in R$ such that
\[
 \gcd\left(m,\text{all entries of }v+\sum_{i=1}^k t_i u_i\right)=\gamma.
\]
Every vector in this affine coset has content divisible by $\gamma$.
Consequently $\gamma$ is its minimum content, both numerically and in the
divisibility order.  A witnessing coefficient tuple can be computed by a
deterministic algorithm polynomial in the binary input length, without
factoring $m$ or enumerating parameter values.
\end{theorem}

A prime-by-prime existence argument would choose parameters modulo each
prime power and invoke the Chinese remainder theorem.  The following
two-vector lemma supplies a direct factor-free implementation instead.

\begin{lemma}[Factor-free content merge]
\label{lem:content-merge}
Let $a,b\in\mathbb Z^\beta$ represent two vectors modulo $m$, and put
\[
 g=\gcd(m,a_1,\ldots,a_\beta,b_1,\ldots,b_\beta),\quad
 M=m/g,\quad d=\gcd(M,a_1/g,\ldots,a_\beta/g).
\]
If $d=1$, take $t=0$.  Otherwise start $D=d$ and repeatedly replace
\begin{equation}
 D\longleftarrow\gcd(M,D^2)
 \label{eq:gcd-saturation}
\end{equation}
until $D$ is unchanged, and take $t=M/D$.  Then
$\gcd(m,a_1+t b_1,\ldots,a_\beta+t b_\beta)=g$.
The saturation loop uses at most $1+\lceil\log_2\log_2 M\rceil$
gcd evaluations when $M\geq2$.
\end{lemma}

\begin{proof}
Write $A=a/g$ and $B=b/g$, so that $\gcd(M,A_1,\ldots,B_\beta)=1$.
For each prime $p$ dividing $M$, write $e=v_p(M)$ and $f=v_p(d)$.
After $j$ saturation updates, the valuation in $D$ is
$\min(e,2^j f)$.  Thus at termination $D$ contains the full prime-power
factors of $M$ supported on the primes of $d$, and no other primes.
The integers $D$ and $Q=M/D$ are coprime.  This description proves the
claimed iteration bound, since $e\leq\log_2 M$; it is a proof of the
algorithm and does not ask the algorithm to find any prime factor.

If $p\mid D$, all entries of $A$ vanish modulo $p$, some entry of $B$ is
nonzero modulo $p$, and $Q$ is a unit modulo $p$.  Hence $A+QB$ has a
nonzero entry modulo $p$.  If $p\mid Q$, some entry of $A$ is nonzero
modulo $p$, while $QB$ vanishes modulo $p$.  In both cases the new vector
has a nonzero entry.  It therefore has content one modulo $M$, proving
the stated content $g$ modulo $m$.  If $d=1$, the vector $A$ already has
content one and no update is needed.
\end{proof}

\begin{proof}[Proof of Theorem~\ref{thm:attained-affine-content}]
Start with $a=v$ and merge successively with $u_1,\ldots,u_k$ using
Lemma~\ref{lem:content-merge}.  After the $j$th merge, the current vector
has content equal to the gcd of $m$ and all entries of
$v,u_1,\ldots,u_j$.  Its coefficient on the previous vector is one, so
the result always remains in the prescribed affine coset.  Reducing
entries and coefficients modulo $m$ after each update bounds every stored
residue by $m$. Products such as $D^2$ have at most twice this bit length;
no unrestricted coefficient growth is hidden in the reduction.  Each merge requires $O(\beta+\log\log(m+2))$ gcd
evaluations and $O(\beta)$ modular arithmetic operations on
$O(\log(m+2))$-bit integers.  Conversely $\gamma$ divides every input
entry and hence every affine combination.  This proves attainment,
minimality, and polynomial binary complexity.
\end{proof}

The saturation is essential: for $m=12$, $a=2$, $b=1$, the unsaturated
choice $t=m/\gcd(m,a)=6$ gives content four, whereas the saturated
choice is $t=3$, giving $a+tb=5$ and content one.  Testing just
$t=0,1$ also fails already for $m=6$, $a=2$, $b=1$.

\subsection{Solving the seam equations with bounded residues}

Classical polynomial algorithms compute integer Smith and Hermite normal
forms together with transformation matrices~\cite{kannan-bachem}.
For the present single-modulus problem, a simpler incremental reduction
suffices and avoids reliance on a particular matrix package.

Maintain a complete parameterization of the equations processed so far,
\[
 z=z_0+B t\pmod m,\qquad t\in R^n.
\]
Initially $z_0=0$ and $B=I$.  Zero or redundant columns may be retained;
injectivity is unnecessary.  For a new equation $a_0 z=b_0$, form
\[
 a=a_0B\pmod m,\qquad \delta=b_0-a_0z_0\pmod m.
\]
If $a=0$, consistency is equivalent to $\delta=0$.  Otherwise swap a
nonzero coefficient into the first position and reduce the row $a$ to
$(d,0,\ldots,0)$ by extended-gcd two-column transformations.  In one step,
with current first coefficient $d_0$ and another coefficient $a_j$,
choose $u d_0+v a_j=g=\gcd(d_0,a_j)$ and use
\[
 \begin{pmatrix}u&-a_j/g\\v&d_0/g\end{pmatrix}.
\]
Its determinant is one.  It changes the parameter coordinates bijectively
over $R$, and the corresponding two columns of $B$ can be updated directly.
All entries of $B$ are reduced modulo $m$ after each operation.

The remaining scalar equation is $d t_1=\delta\pmod m$.  Put
$g=\gcd(d,m)$ and $Q=m/g$.  There is a solution exactly when $g\mid\delta$.
If so, choose
\[
 t_1^{(0)}=(\delta/g)(d/g)^{-1}\pmod Q,
\]
using $t_1^{(0)}=0$ when $Q=1$.  Update $z_0$ by
$t_1^{(0)}$ times the first column of $B$, and multiply that column by
$Q$.  This gives exactly all solutions of the enlarged system.  The
parameter dimension never grows beyond $n$.

\begin{proposition}[Exact feasible-assignment count]
Starting with $N=m^n$, divide $N$ by $Q=m/\gcd(m,a_1,\ldots,a_n)$
at every feasible nonzero scalar restriction above.  The final value is
$|S|$, the number of distinct assignments satisfying the seam equations.
\end{proposition}

\begin{proof}
The image of the new scalar functional on the previous solution module
is the subgroup of $R$ of size $Q$.  Every nonempty fibre of a homomorphism
between finite groups has the same cardinality.  The new affine restriction
therefore keeps exactly the fraction $1/Q$ of the previous solutions.
Although the maintained parameterization can be redundant, all of its
fibres over the previous solution module also have equal size, so the
same fraction applies to distinct assignments.  A zero consistent row
has $Q=1$ and makes no change.
\end{proof}

After all seams have been solved, set
\[
 v=c+Hz_0,\qquad U=HB\pmod m.
\]
Theorem~\ref{thm:attained-affine-content} gives the exact minimum of
\eqref{eq:affine-component-objective} and a witness $z_*=z_0+Bt_*$.
In particular, a connected member exists if and only if $S$ is nonempty
and
\begin{equation}
 \gcd\bigl(m,\text{all entries of }v,U\bigr)=1.
 \label{eq:affine-connected-criterion}
\end{equation}
The constrained solver and the witness construction are both deterministic
and polynomial in $n,q,\beta$ and the binary coefficient lengths.
For example, with classical arithmetic a conservative polynomial bound
follows from $O(qn^2)$ modular matrix operations, $O(qn)$ extended gcds,
$O(\beta n^2)$ operations for the image, and
$O(n(\beta+\log\log(m+2)))$ gcds for the optimization. The final
witness multiplication adds $O(n^2)$ modular operations. The auxiliary
count $N\le m^n$ has $O(n\log(m+2))$ bits, and its at most $q$ exact
divisions also have polynomial bit cost. This is not a
claim of polynomial running time in the number of tetrahedra of an
arbitrary knot exterior: constructing a suitable complete family remains
a separate geometric problem.

